% vim:ft=tex:
%
\RequirePackage[l2tabu, orthodox]{nag}
\documentclass[12pt]{llncs}
\pagestyle{headings}


%%%%%%%%%%%%%%%%%%%%%%%%%%%%%%%%%%%%%%%
% FONTS
%%%%%%%%%%%%%%%%%%%%%%%%%%%%%%%%%%%%%%%

\usepackage{lmodern}        % Improved version of the original Computer Modern font

\usepackage[T1]{fontenc}    % Determines font encoding of the output. Font packages have to be included before this line.
\usepackage[utf8]{inputenc} % Determines encoding of the input. All input files have to use UTF8 encoding.
\usepackage[english]{babel}


%%%%%%%%%%%%%%%%%%%%%%%%%%%%%%%%%%%%%%%
% PACKAGES
%%%%%%%%%%%%%%%%%%%%%%%%%%%%%%%%%%%%%%%

\usepackage{amsmath}
\usepackage{amssymb}
% \usepackage{thmtools}
% \usepackage{mathtools}
% \usepackage{prftree}
% \usepackage{cancel}
% \usepackage{keycommand}
\usepackage{listings}
% \usepackage{lstautogobble}
\usepackage{microtype}  % Small-scale typographic enhancements.
\usepackage[shortlabels]{enumitem}
\usepackage[hidelinks]{hyperref}  % optional argument [hidelinks] hides colored boxes around clickable areas
\usepackage{prftree}
\usepackage{mathtools}
\usepackage{booktabs}   % Improves the typesettings of tables.
\usepackage{tabularx}
\usepackage{multirow}   % Allows table elements to span several rows.
% \usepackage{pdfpages}  % to include pages from existing pdf files with \includepdf[pages={1,2}]{./path/to/file.pdf}
\usepackage{xcolor}
\usepackage{todonotes}  % Add optional argument [disable] to hide TODOs, or [obeyFinal] to hide only in final mode
\usepackage{mdframed}

\usepackage{tikz}
\usetikzlibrary{arrows, arrows.meta, backgrounds, calc, decorations.markings, positioning, shapes.geometric}


%%%%%%%%%%%%%%%%%%%%%%%%%%%%%%%%%%%%%%%
% SETTINGS
%%%%%%%%%%%%%%%%%%%%%%%%%%%%%%%%%%%%%%%

% \setlist{nolistsep}  % less vertical spacing for itemize/enumerate environments


%%%%%%%%%%%%%%%%%%%%%%%%%%%%%%%%%%%%%%%
% MACROS
%%%%%%%%%%%%%%%%%%%%%%%%%%%%%%%%%%%%%%%

\makeatletter
% \def\UrlFont{\rmfamily}
% \def\orcidID#1{\smash{\href{http://orcid.org/#1}{\protect\raisebox{-1.25pt}{\protect\includegraphics{orcid_color.eps}}}}}
% The \smash breaks the hyperlink... the clickable area becomes tiny.
\def\orcidID#1{{\href{http://orcid.org/#1}{\protect\raisebox{-1.25pt}{\protect\includegraphics{ORCID_Color.eps}}}}}
\makeatother

\newcommand{\todoi}[1]{\todo[inline,caption={}]{#1}}    % caption={} allows itemize etc. inside todoi, see https://tex.stackexchange.com/a/54068

% This is the provability analogue to \models (note that \vdash is smaller)
\DeclareRobustCommand\proves{\mathrel{|}\joinrel\mkern-.5mu\mathrel{-}}

\newcommand{\limpl}{\rightarrow}    % logical implication
\newcommand{\Limpl}{\Rightarrow}    % logical implication (variant)
\newcommand{\liff}{\leftrightarrow} % logical equivalence
\newcommand{\Liff}{\Leftrightarrow} % logical equivalence (variant)

\newcommand{\Land}{\bigwedge}
\newcommand{\Lor}{\bigvee}

\newcommand{\union}{\cup}
\newcommand{\Union}{\bigcup}
\newcommand{\intersect}{\cap}
\newcommand{\Intersect}{\bigcap}

\newcommand{\eql}{\simeq}
\newcommand{\neql}{\not\eql}

% \newcommand{\nat}{\mathbb{N}}
% \newcommand{\int}{\mathbb{Z}}

\newcommand{\naf}{{\sim}}           % negation as failure, default negation
\newcommand{\caF}{\mathcal{F}}
\newcommand{\caM}{\mathcal{M}}
\newcommand{\caT}{\mathcal{T}}

\newcommand{\vampire}{\textsc{Vampire}}


%%%%%%%%%%%%%%%%%%%%%%%%%%%%%%%%%%%%%%%
% METADATA
%%%%%%%%%%%%%%%%%%%%%%%%%%%%%%%%%%%%%%%

\title{Automated Reasoning in Generating Exam Sheets for Automated Deduction}
\author{Petra Hozzov\'a\and
Laura Kov\'acs \and
Jakob Rath}

\institute{
    TU Wien, Austria
    % \and
    % Princeton University, Princeton NJ 08544, USA \and
    % Springer Heidelberg, Tiergartenstr. 17, 69121 Heidelberg, Germany
    % \email{lncs@springer.com}\\
    % \url{http://www.springer.com/gp/computer-science/lncs} \and
    % ABC Institute, Rupert-Karls-University Heidelberg, Heidelberg, Germany\\
    % \email{\{abc,lncs\}@uni-heidelberg.de}
}

% Set PDF document properties
\hypersetup{
    pdfpagemode     = UseNone,                  % Don't show bookmarks when opening the pdf file, see also https://latex.org/forum/viewtopic.php?p=70568#p70568
    pdfpagelayout   = OneColumn,                % How the document is shown in PDF viewers (optional). Values: SinglePage, OneColumn, TwoColumnLeft, TwoColumnRight, TwoPageLeft, TwoPageRight
    % pdfauthor       = {\authorname},          % The author's name in the document properties (optional).
    % pdftitle        = {\thesistitle},         % The document's title in the document properties (optional).
    % pdfsubject      = {subject},              % The document's subject in the document properties (optional).
    % pdfkeywords     = {keyword, another keyword}, % The document's keywords in the document properties (optional).
    colorlinks,
    % linkbordercolor = {Melon},                % The color of the borders of boxes around crosslinks (optional).
    % linkcolor={black},
    % citecolor={black},
    % urlcolor={black},
}

\lstset{%
    basicstyle=\ttfamily,%
    % keywordstyle=\bfseries\color{blue},%
    columns=flexible,%
    % numberstyle=\tiny,%
    % numbersep=5pt,%
    % numbers=left,%
}


%%%%%%%%%%%%%%%%%%%%%%%%%%%%%%%%%%%%%%%
\begin{document}
\maketitle





\begin{abstract}
Amid the COVID-19 pandemic, distance teaching became the default in world-wide higher education,
urging teachers and researchers to revise course materials into a
more accessible online content for a diverse audience. Probably one of the hardest challenges in this new form of education
came with online assessments of course performance, especially organizing and grading online written exams.
In this short paper we describe the online assesment
setting  we organized for our master's level course ``Automated
Deduction'' at TU Wien.
The algorithmic and rigorous logical reasoning developed within our
course calls for exam sheets focused on problem solving and deductive
proofs; as such exam sheets using test grids are not a viable solution
for written exams within our course.
We report on the automated reasoning framework we developed for generating individual online exam sheets for students enrolled in the course.
We believe that the toolchain of automated reasoning tools we have developed for holding online written exams could be beneficial not only for other distance learning platforms, but also to researchers in automated reasoning by providing our community with a large set of randomly generated benchmarks in SAT/SMT solving and first-order theorem proving.
\end{abstract}





\section{Motivation}

online exams due to pandemic.
we want to avoid collaboration between students during exam (or make it at least a bit harder),
so each student gets their own exam sheet.
etc…





\section{todo}

Challenge:
problem instances should be different but of similar difficulty to make sure the exam is fair to students.





\section{Method 1: Varying Templates and Fixed Patterns}


\subsection{Satisfiability Modulo Theories (SMT)}

relatively easy: provided template, do small random perturbations that don't change the solution


\subsection{Ground Superposition}

todo




\section{Method 2: Full Random Generation}

more sophisticated: full random generation, filter out "too hard"/"too easy" instances.
Note that we don't need very efficient implementation of filters since the instances are very small.
so we can use naive satisfiability tests or model counting.



\subsection{Boolean Satisfiability (SAT)}

% \fbox{\parbox{\textwidth-10pt}{%
\begin{mdframed}
\begin{example}[SAT]
    Hello
Consider the formula:
\[
% Random number generator seed: Just 123
% Index: 36
% Satisfiable? True
% Valid? False
% Polarities: fromList [(p,Both),(q,Both),(r,Pos)]
% Models: [{¬p,q,r}, {¬p,¬q,r}, {¬p,¬q,¬r}]
( r \land \lnot ( q \rightarrow p ) ) \lor ( q \leftrightarrow \lnot ( p \rightarrow q ) )
\]
\begin{itemize}
\item[(a)]
Which atoms are pure in the above formula?
\item[(b)] Compute a clausal normal form $C$ of the above formula by
  applying the CNF transformation algorithm with naming and
  optimization based on polarities of subformulas;
\item[(c)] Decide the satisfiability of the computed CNF formula $C$
  by applying  the DPLL method to $C$. If $C$ is
satisfiable,  give an interpretation which
satisfies it.
    \qed
\end{itemize}
\end{example}
\end{mdframed}
% }}

describe filters, and why they were chosen (aim for a challenging problem, but still solvable by hand).

problems:
when restricting too much, the resulting formulas may end up boring.
e.g. SAT formula with exactly one model, or no model


\subsection{Non-Ground Superposition and Redundancy}

todo


\section{Implementation}

Various programs/scripts tied together by a control script.
Running the control script generates all problems and puts them together, resulting in $n$ different exams in PDF format.


Latex template:
leave holes for formulas/signatures that will be inserted using the command ``\textbackslash{}input''.


Problem templates: (todo: describe impl of smt and groundsup)


For the full random generation of the SAT and redundancy problems,
we were initially inspired by Haskell's QuickCheck library.
\todo{cite?
Claessen, Koen and Hughes, John (2000). "QuickCheck: A Lightweight Tool for Random Testing of Haskell Programs" (PDF). Proceedings of the International Conference on Functional Programming (ICFP), ACM SIGPLAN.
}
We used QuickCheck for the first prototype, but realized that we need backtracking in the generator.

Added mtl-style typeclass 'MonadChoose' with a primitive operation 'choose' for choosing an element from a list of choices.
Our generator implementations are generic over the monad, constrained by this typeclass.

We used two concrete implementations to evaluate generators:
1. 'RandomChoiceT', a monad transformer that implements 'choose' as uniform random choice with backtracking support
    (can be thought of as 'ListT' where 'choose' is like the normal monadic bind but shuffles the list with a random permutation first).
    This is what we actually use to generate random exams.
2. a standard list monad (or similar) to enumerate the sample space.
    This second evaluation method lets us ensure that the sample space is sufficiently large.

% Thus, re-implemented nondeterministic monad with an additional primitive operation for uniform random choice of an element from a list.
% We used mtl-style typeclass for the generator implementations;
% this allows us to use the same generator in two ways.
% 1. get a random element (by evaluating in a 
% 2. list all possible elements (by evaluating in the built-in list monad, or a similar one)
% The second evaluation method lets us ensure that the sample space is sufficiently large.


\section{Conclusion}



% \bibliographystyle{splncs04}
% \bibliography{references}
% % \bibliography{/Users/jakob/Documents/Library/library.bib}

\end{document}
